% sec-01-abstract.tex - Abstract and Keywords (full)
\begin{abstract}
\noindent
A Phase 1 oncology trial advances only when a sequence of large regulatory and
clinical documents is authored, reviewed, and approved, and patients already
enrolled cannot wait while that paperwork is assembled. This paper shows how a
repository based large language model (LLM), driven by a single master prompt that
first writes and then executes a schedule of sub-prompts, can hasten the entire
Phase 1 process by generating every relevant large document through an auditable
mermaid to draft to full to final pipeline. The method is demonstrated on a real
pancreatic ductal adenocarcinoma (PDAC) program: an on-premises LLM directed
robotic Whipple with perioperative daraxonrasib (RMC-6236), whose Phase 1 and
Phase 2 protocols were themselves single-prompt multi-file builds. We map the six
document targets where faster, validated authoring yields the greatest schedule
value (the initial IND and IRB package; protocol amendments with synchronized
consent; cohort-review packages; the complete clinical-hold response; the
Phase 2-to-3 briefing package; and the pivotal CSR and NDA/BLA modules), and we
show that only the administrative and preparation time bucket compresses, while
clinical follow-up and fixed regulatory review clocks do not. New document
iterations complete in one to four days. Utility is verified, not asserted:
twenty-four colored Mermaid figures ground the prose and reproduce identically in
LaTeX, paper repository, directory, and file names match the repository, and every
file is committed to GitHub in real time for monitorability. For enrolled patients
who cannot wait, probable benefit exceeds probable risk. Paper v1.0; repository
v4.2.0.
\end{abstract}

\keywords{Pancreatic ductal adenocarcinoma; Phase 1 clinical trial; large language
model; document generation; repository based LLM; sub-prompt schedule; mermaid to
LaTeX; clinical trial acceleration; daraxonrasib; verification before generation.}
